\begin{table*}[t]
\centering
\scriptsize
\renewcommand{\arraystretch}{0.9}
\begin{tabularx}{\linewidth}{>{\raggedright\arraybackslash}p{0.19\linewidth}>{\raggedright\arraybackslash}p{0.18\linewidth}>{\raggedright\arraybackslash}p{0.18\linewidth}>{\raggedright\arraybackslash}p{0.20\linewidth}X}
\toprule
\textbf{Field} & \textbf{coder~1 vs.\ final} & \textbf{coder~2 vs.\ final} & \textbf{coder~1 vs.\ coder~2} & \textbf{Final adjudication effect} \\
\midrule
Primary failure type & 0.784 / 0.743 & 0.899 / 0.874 & 0.764 / 0.719 & Final labels resolve 125 rows that changed from the candidate expansion on at least one field. \\
Method family & 0.878 / 0.818 & 0.946 / 0.915 & 0.838 / 0.762 & Broad candidate method labels were corrected toward coder-supported labels. \\
Intervention point & 0.838 / 0.807 & 0.885 / 0.863 & 0.770 / 0.724 & Pipeline-location disagreements are resolved before final Map-B synthesis. \\
Scope / inclusion & 0.953 / 0.905 & 0.932 / 0.865 & 0.885 / 0.769 & Final Map B retains 74 additions and excludes 74 boundary candidates. \\
\bottomrule
\end{tabularx}
\caption{Final-label agreement for the 148 Map-B candidate additions. Cells report agreement / Cohen's $\kappa$ against the final adjudicated label, not the pre-adjudication candidate label. This resolves the apparent instability of the candidate map: coder~1 and coder~2 agree strongly with the final labels, and excluded boundary rows are kept in the release but not in canonical Map B.}
\label{tab:map-b-second-coding-c02-c03}
\end{table*}

